\documentclass[12pt,a4paper]{article}
\usepackage[a4paper,left=22mm,right=22mm,top=22mm,bottom=23mm]{geometry}
\usepackage{fontspec}
\setmainfont{DejaVu Serif}
\usepackage[russian]{babel}
\usepackage{amsmath,amssymb,mathtools}
\usepackage{array,booktabs,longtable}
\usepackage{xcolor}
\usepackage{tikz}
\usepackage[unicode,colorlinks=true,linkcolor=blue!55!black]{hyperref}
\setlength{\parindent}{0pt}
\setlength{\parskip}{0.48em}
\renewcommand{\arraystretch}{1.35}
\setcounter{secnumdepth}{0}
\hypersetup{pdftitle={Ревью домашней работы},pdfauthor={Лёвин Артём Александрович}}
\begin{document}
\begin{titlepage}
\centering
\vspace*{\fill}
{\Huge\bfseries Ревью домашней работы\par}
\vspace{1.7cm}
{\Large Автор: Лёвин Артём Александрович\par}
\vspace{0.75cm}
{\large 9 октября 2026 года\par}
\vspace*{\fill}
\begin{tikzpicture}
  \draw[blue!40!black,line width=0.7pt]
    (-4.0,0) .. controls (-2.2,0.65) and (-0.9,-0.45) .. (0,0)
             .. controls (1.2,0.55) and (2.4,-0.45) .. (4.0,0);
\end{tikzpicture}
\vspace{0.8cm}
\end{titlepage}

\newpage
\section*{Сводная таблица ошибок}
В таблицу внесены только задания, для которых требуется исправление.

\small
\begin{tabular}{@{}p{1.15cm}p{3.35cm}p{5.1cm}p{2.05cm}p{2.05cm}@{}}
\toprule
\textbf{№} & \textbf{Тема} & \textbf{Суть ошибки} & \textbf{Ваш ответ} & \textbf{Правильный ответ} \\
\midrule
\hyperref[task:5]{\textbf{5}} & Путь за отдельную секунду равноускоренного движения & В формулах пути за первые четыре и первые пять секунд использована одна и та же конечная скорость. & $36\ \text{м/с}$ & $20\ \text{м/с}$ \\
\bottomrule
\end{tabular}
\normalsize

\newpage
\thispagestyle{empty}
\vspace*{\fill}
\begin{center}
{\Huge\bfseries Разбор ошибок}
\end{center}
\vspace*{\fill}

\newpage
\phantomsection\label{task:5}
\section*{Задание № 5}

\textbf{Текст задания.} Тело, двигаясь из состояния покоя, за пятую секунду прошло $18\ \text{м}$. Какую скорость оно имеет в конце пятой секунды?

\textbf{Уточнение.} В самом условии не сказано прямо, что ускорение постоянно. Ниже, в соответствии с темой данного набора задач, рассматривается \emph{равноускоренное движение}. Без такого предположения однозначно определить конечную скорость по приведённым данным нельзя.

\textbf{Ваше решение.} Вы справедливо начали с того, что путь за пятую секунду необходимо получить как разность путей за первые пять и первые четыре секунды. Далее была записана формула пути через среднюю скорость:
\[
S=\frac{(v+v_0)t}{2},\qquad v_0=0.
\]
Затем Вы записали следующие преобразования:
\[
\begin{aligned}
18&=S(5)-S(4),\\
18&=\frac{v\cdot5}{2}-\frac{v\cdot4}{2},\\
18&=2{,}5v-2v,\\
18&=0{,}5v,\\
v&=36\ \text{м/с}.
\end{aligned}
\]

\textbf{Ваш ответ.} $36\ \text{м/с}$.

\textcolor{red}{\textbf{Ошибка:}} при вычитании путей за два различных промежутка времени для обоих промежутков использована одна и та же конечная скорость $v$.

\textbf{Где именно возникает ошибка.}

Последний верный шаг:
\[
\boxed{S_{\text{пятая секунда}}=S(5)-S(4)=18\ \text{м}.}
\]
Это правильное разбиение задачи: путь за пятую секунду равен разности пути, пройденного к концу пятой секунды, и пути, пройденного к концу четвёртой секунды.

Первый неверный переход:
\[
\boxed{18=\frac{v\cdot5}{2}-\frac{v\cdot4}{2}.}
\]
В первом слагаемом символ $v$ обозначает скорость \emph{в конце пятой секунды}. Во втором слагаемом вместо него должна стоять скорость \emph{в конце четвёртой секунды}. Эти скорости не равны при ненулевом постоянном ускорении.

\textbf{Почему этот шаг неверен.}

Формула
\[
S(t)=\frac{v_0+v(t)}{2}\,t
\]
для равноускоренного движения использует начальную скорость $v_0$ и скорость $v(t)$ \emph{в конце именно того интервала}, для которого считается путь. Значит, правильно писать:
\[
S(5)=\frac{v_0+v_5}{2}\cdot5,
\qquad
S(4)=\frac{v_0+v_4}{2}\cdot4,
\]
где $v_5$ --- скорость к концу пятой секунды, а $v_4$ --- скорость к концу четвёртой секунды. При движении из состояния покоя с постоянным ускорением $a$:
\[
v_4=4a,\qquad v_5=5a.
\]
Поэтому без дополнительного обоснования нельзя заменить $v_4$ на $v_5$. Полученное значение $36\ \text{м/с}$ стало следствием именно этой подмены, а не ошибки при последнем арифметическом делении.

\textcolor{blue!60!black}{\textbf{Ключевая идея:}} для вычисления пути \emph{за пятую секунду} нужно вычесть путь за первые четыре секунды из пути за первые пять секунд. Каждый из этих двух путей вычисляется со своим временем и со своей соответствующей этому времени скоростью. Ещё проще воспользоваться формулой координаты (перемещения) при постоянном ускорении.

\textbf{Исправление от последнего правильного шага.}

Сохраняем верное равенство
\[
S_{\text{пятая секунда}}=S(5)-S(4)=18\ \text{м}.
\]
Для движения из состояния покоя с постоянным ускорением
\[
S(t)=v_0t+\frac{at^2}{2}=\frac{at^2}{2},
\qquad v_0=0.
\]
Подставим отдельно время $5\ \text{с}$ и время $4\ \text{с}$:
\[
\begin{aligned}
18&=\frac{a\cdot5^2}{2}-\frac{a\cdot4^2}{2}\\
&=\frac{25a-16a}{2}\\
&=\frac{9a}{2}=4{,}5a.
\end{aligned}
\]
Отсюда
\[
a=\frac{18}{4{,}5}=4\ \text{м/с}^2.
\]
В конце пятой секунды скорость равна
\[
v_5=v_0+a\cdot5=0+4\cdot5
=\boxed{20\ \text{м/с}}.
\]

\textbf{Полное правильное решение.}

\emph{Шаг 1.} Обозначим ускорение тела через $a$. Поскольку тело начинает движение из состояния покоя, его начальная скорость равна нулю. При постоянном ускорении путь с момента начала движения до момента $t$ составляет
\[
S(t)=\frac{at^2}{2}.
\]

\emph{Шаг 2.} Интервал \emph{пятой секунды} начинается при $t=4\ \text{с}$ и заканчивается при $t=5\ \text{с}$. Поэтому пройденные за эти две границы пути:
\[
S(4)=\frac{16a}{2}=8a,
\qquad
S(5)=\frac{25a}{2}=12{,}5a.
\]

\emph{Шаг 3.} Вычислим разность и приравняем её данному расстоянию:
\[
S(5)-S(4)=12{,}5a-8a=4{,}5a=18.
\]
Ускорение тела составляет $4\ \text{м/с}^2$.

\emph{Шаг 4.} Скорость при равноускоренном движении находится по формуле
\[
v(t)=v_0+at.
\]
При $t=5\ \text{с}$ получаем
\[
v(5)=0+4\cdot5=20\ \text{м/с}.
\]

\textbf{Проверка.}

К концу четвёртой секунды скорость равна $16\ \text{м/с}$. К концу пятой секунды она равна $20\ \text{м/с}$. Средняя скорость за пятую секунду составляет
\[
v_{\text{ср}}=\frac{16+20}{2}=18\ \text{м/с}.
\]
За одну секунду тело проходит
\[
S_{\text{пятая секунда}}=18\ \text{м/с}\cdot1\ \text{с}=18\ \text{м}.
\]
Это точно совпадает с условием, поэтому найденное значение скорости согласуется с исходными данными.

\textcolor{teal!60!black}{\textbf{Итог:}} при равноускоренном движении из состояния покоя скорость в конце пятой секунды равна
\[
\boxed{v_5=20\ \text{м/с}}.
\]

\subsection*{Задача на закрепление}
Тело начинает равноускоренное движение из состояния покоя. За седьмую секунду оно проходит $39\ \text{м}$. Найдите скорость тела в конце седьмой секунды.

\vfill
\begin{center}
\small Лёвин Артём Александрович эксклюзивно для Ивана Петраченкова
\end{center}
\end{document}
